% Pursuing Stacks (cotes 134-2 to 134-8): archive cote 134-6, site page 89 (PDF page 90), typescript page 407.
% typescript. Draft compared: ps5.tex:5999-6080.
% AI-transcribed by Claude Opus 5.5 (claude-opus-5-5, adaptive thinking, effort high), Message Batches API, batch ps3 (msgbatch_01XdFmuVVBAcgc7EFxdtCqsj), 2026-11-12; not yet checked by eye.
\scanpage{134-6/89}
\mspage{407}

14.8.\hand{ and 15.8.}

106. \hand{E) \uline{A formulaire around the basic operations} $*$ \uline{and} Hom.}

I would like to dwell a little more still on the duality formalism
weaving around formula (39) (p.395), stating that for two \inserted{given} small categorie
% NOTE: right edge of the scan cuts the end of this line after "categorie"
A and B dual to each other
\[ B=A^{o} \ , \ A=B^{o} \quad , \]
and any \underline{U}-category \hand{$\mathcal{M}$} stable under both types of small limits, \hand{$\mathcal{M}$}-valued
\uline{cosheaves} on A\^{} \over{\unsure{may}}{may} be interpreted as \hand{$\mathcal{M}$}-valued \uline{sheaves} on the dual topos
B\^{}. This identification preserves variance (i.e. (39) is an \uline{equivalence}
of categories, not an antiequivalence :
\[ (59) \qquad A^{\wedge\hand{\mathcal{M}}} = \underline{\mathrm{Hom}}_{\hand{!}}(A^{\wedge},\hand{\mathcal{M}}) \hand{\xrightarrow{\approx}} \underline{\mathrm{Hom}}^{\hand{!}}(B^{\hand{\unsure{\wedge}o}},\hand{\mathcal{M}}) \ = B^{\wedge}_{\hand{\mathcal{M}}} \qquad ). \]
It should not be confused with the tautological interpretation of
\hand{$\mathcal{M}$}-valued cosheaves on A\^{} as \hand{$\mathcal{M}^{o}$}-valued sheaves on the \uline{same} topos, \xout{\illegible{1}}\inserted{an iden-}
\deleted{\unsure{pending}}\inserted{tification} \uline{reversing} \uline{variance}, as expressed by the canonical \inserted[typed]{\xout{anti}} \xout{isomorphism}
antiequivalence between the corresponding categories -- an antiisomor\xout{phis}
phism even (reflecting it's tautological nature):
\margin{$\hat{A}_{\mathcal{M}^{o}}$}
\[ (60) \qquad \underline{\mathrm{Hom}}_{\hand{!}}(A^{\wedge},\hand{\mathcal{M}})^{o} \hand{\xrightarrow{\sim}} \underline{\mathrm{Hom}}^{\hand{!}}(A^{\wedge o},\hand{\mathcal{M}^{o}}) \qquad \text{i.e.} \quad (A^{\wedge\hand{\mathcal{M}}})^{o} \hand{\xrightarrow{\approx}} \deleted{\illegible{1}} \hand{\hat{A}_{\mathcal{M}^{o}}}. \]
In the latter formula, the basic topos A\^{} remains the same in both sides,
it is the category of values which changes from \hand{$\mathcal{M}$} to the dual one \hand{$\mathcal{M}^{o}$} ,
whereas in formula (59) = (39), it is the opposite. In terms of the tau-%
tological formula (60) (a particular case of formula (6) p.385), the
not-so-tautological formula \unsure{relying} cosheaves and sheaves can be reformu-%
lated as a formula in terms of sheaves only (due to our preference for
sheaves rather than cosheaves ...):
\[ (61) \qquad \underline{\mathrm{Hom}}^{\hand{!}}(A^{\wedge},\hand{\mathcal{M}})^{o} \hand{\approx} \underline{\mathrm{Hom}}^{\hand{!}}(B^{\wedge},\hand{\mathcal{M}^{o}}) \qquad \text{i.e.} \quad (A^{\wedge}_{\hand{\mathcal{M}}})^{o} \hand{\approx} \deleted{(}B^{\wedge}\deleted{\illegible{1}}{}_{\hand{\mathcal{M}^{o}}} \ , \]
namely \inserted[typed]{\hand{$\mathcal{M}$}-valued\hand{)}} sheaves on the topos A\^{} can be interpreted as sheaves on the dual
topos with values in the dual category \hand{$\mathcal{M}^{o}$} , this interpretation \uline{rever-}%
\uline{sing} \uline{variances}. In the homology and cohomology formalism which is to
follow, due to \xout{my} habits of long standing, I prefer systematically to
take as coefficients \uline{sheaves} rather than cosheaves -- hence \xout{\unsure{interprete}}
rule out cosheaves in favor of sheaves via (60). From this point of view,
the relevant basic duality statement is (61) rather than (59).

On the \uline{cosheaves} side, yesterday's diagram (58) of equivalences
gives a fourfold description of cosheaves on the topos A\^{} with values in
an \uline{abelian} category \hand{$\mathcal{M}$} stable under small direct limits. We could still
enlarge this diagram, by \xout{adding} including in it a fifth category equiva-%

\margin{89}
